\section{Input file\label{sec:Input-file}}

In order to run JETSPIN simulations an input file has to be prepared,
which is a free-format and case-insensitive. The input file has to
be named \texttt{input.dat}, and it contains the selection of the
model system, integration scheme directives, specification of various
parameters for the model, and output directives. The input file does
not requires a specific order of key directives, and it is read by
the input parsing module. Every line is treated as a command sentence
(record). Records beginning with the symbol \# (commented) and blank
lines are not processed, and may be added to aid readability. Each
record is read in words (directives and additional keywords and numbers),
which are recognized as such by separation by one or more space characters.

As in the example given in \texttt{input test}, the last record is
a \texttt{finish} directive, which marks the end of the input data.
Before the \texttt{finish} directive, a wide list of directives may
be inserted (see Tab \ref{tab:inputlist}). The key \texttt{systype}
should be used to set the nanofiber model. In JETSPIN two models are
available: 1) the one dimensional model similar to the model of Cap
\ref{chap:JETSPIN-Model-and} but assuming the nanofiber to be straight
along the $\vec{\textbf{x}}$ axis, and, therefore, neglecting the
surface tension force $\vec{\textbf{f}}_{st}$; 2) the three dimensional
model described in Cap \ref{chap:JETSPIN-Model-and}. Internally these
options are handled by the integer variable \texttt{systype}, which
assumes the values explained in Tab \ref{sec:Input-file}. Further
details of the 1-D model can be found in Refs \cite{pontrelli2014electrospinning,carroll2011discretized}.
A series of variables are mandatory and have to be defined. For example,
\texttt{timestep}, \texttt{final time}, \texttt{initial length}, etc.
(see underlined directives in Tab \ref{tab:inputlist}). A missed
definition of any mandatory variable will call an error banner on
the terminal. The user can use two possible ways to define the initial
geometry of a nanofiber both given the mandatory directive \texttt{initial
length}: 1) the discretization step length of nanofiber is specified
by the directive \texttt{resolution}, which causes automatically the
setting of the jet segments number (the number of segment in which
the jet is discretised); 2) the jet segments number is declared by
the directive \texttt{points}, while the value of the discretization
step length is automatically set by the program (as in Example \textbf{Test
Case 3} in Sec \ref{sec:Test-Case-3:}). The directive \texttt{cutoff}
indicates the length of the upper and lower proximal jet sections,
which interact via Coulomb force on any bead. It is worth stressing
that the length value is set at the nozzle, so that the effective
cutoff increases along the simulation as the nanofiber is stretched.
The directive \texttt{restart yes} can be used to restart a JETSPIN
job (for further information see Section \ref{sec:Restarting-JETSPIN}).

The user should pay special attention in choosing the time integration
step given by the directive \texttt{timestep}, whose detailed considerations
are provided in Ref \cite{lauricella2015jetspin}. The reader is referred
to Tab \ref{tab:inputlist} for a complete listing of all directives
defining the electrospinning setup parameters. Not all these quantities
are mandatory, but the user is informed that whenever a quantity is
missed, it is usually assumed equal to zero by default (exceptions
are stressed in Tab \ref{tab:inputlist}). \textbf{Note all the quantities
have to be expressed in centimeter\textendash gram\textendash second
unit system} (e.g. charge in statcoulomb, electric potential in statvolt,
etc.).

The \texttt{job time} and \texttt{close time} directives are required
to ensure a controlled close down procedure when a job runs out of
time. The time specified by the job time directive indicates the total
time allowed for the job. (This must obviously be set equal to the
time specified to the operating system when the job is submitted.)
The \texttt{close time} directive represents the time JETSPIN will
require to write and close all the data files at the end of processing.
This means the effective processing time limit is equal to the job
time minus the close time.

\begin{longtable}{cc}
\caption{Here, we report the list of directives available in JETSPIN. Note
$i$, $f$, and $s$ denote an integer number, a floating point number,
and a string, respectively. The underlined directives are mandatory.
The default value is zero (exceptions are stressed in parentheses).}
\tabularnewline
\endlastfoot
\hline
\textbf{directive:}  & \textbf{meaning:}\tabularnewline
\hline
\hline
airdrag yes  & active the inclusion of the airdrag force in the model (default: no)\tabularnewline
 & (note that if yes the directive ``integrator 4'' is mandatory)\tabularnewline
airdrag airdensity $f$  & mass density of the air surrounding the nanofiber (default: 0)\tabularnewline
airdrag airviscosity $f$  & kinematic viscosity of the air surrounding the nanofiber (default:
1)\tabularnewline
airdrag airvelocity $f$  & velocity of the air with respect to the \textit{z} axis (default:
0)\tabularnewline
airdrag amplitude $f$  & amplitude of the randomic force rescaled by $\gamma$ (default: 1)\tabularnewline
 & (note that the dissipation coefficient $\gamma$ is automatically\tabularnewline
 & computed by JETSPIN using the input data by Eq. \ref{eq:fric_dissip})\tabularnewline
close time $f$ & Set job closure time to $f$ seconds\tabularnewline
\uline{collector distance} $f$  & distance of collector along the x axis from the nozzle\tabularnewline
 & \quad{}(note the nozzle is assumed in the origin)\tabularnewline
cutoff $f$  & length of the proximal jet sections interacting by Coulomb \tabularnewline
 & \quad{}force (default: equal to the collector distance)\tabularnewline
\uline{density mass} $f$  & mass density of the nanofiber\tabularnewline
\uline{density charge} $f$  & electric volume charge density of the nanofiber\tabularnewline
dynam refin yes  & active the dynamic refinement during the simultion\tabularnewline
dynam refin threshold $f$  & set the threshold beyond which the refinement is applied\tabularnewline
 & (recommended $\geq20\times$ the discretization resolution; below that\tabularnewline
 & a warning is issued -- see Sec.~\ref{subsec:Refinement-robustness})\tabularnewline
dynam refin every $f$  & the refinement is applied every $f$ seconds\tabularnewline
dynam refin start $f$  & the refinement is applied only if the simulation time is greater \tabularnewline
 & than or equal to $f$ seconds\tabularnewline
dynam refin anchor $f$ & during the refinement procedure all the anchor beads inserted\tabularnewline
 & at distance $f$ preserve their positions\tabularnewline
 & (they are manteined as knots of both the old and new meshes)\tabularnewline
 & (default, if omitted: $5\times$ the discretization resolution;\tabularnewline
 & values below $1\times$ resolution are raised to that same $5\times$ floor)\tabularnewline
dynam refin capacity $i$ & number of beads by which array capacity grows/is reserved\tabularnewline
 & at an accepted refinement event (default: 100; must be positive)\tabularnewline
evaporation yes  & activate solvent evaporation and concentration-dependent rheology (default: no)\tabularnewline
evaporation polymer frac $f$  & initial polymer mass fraction $c_{p,0}$\tabularnewline
evaporation airviscosity $f$  & kinematic viscosity of air $\nu_a$ in $\mathrm{cm}^{2}/\mathrm{s}$\tabularnewline
evaporation temperature $f$  & temperature in kelvin used for the saturation-vapour correlation\tabularnewline
evaporation umidity $f$  & relative humidity $RH$ in the interval $[0,1]$\tabularnewline
 & (historical spelling retained for input-file compatibility)\tabularnewline
evaporation diffusivity $f$  & solvent-vapour diffusion coefficient $D_a$ in $\mathrm{cm}^{2}/\mathrm{s}$\tabularnewline
evaporation relumidity $f$  & optional saturation-vapour ratio $c_{s,\mathrm{eq}}(T)$\tabularnewline
 & (if omitted it is evaluated from \texttt{evaporation temperature})\tabularnewline
evaporation bconstant $f$  & parameter $B$ in the concentration-dependent viscosity law\tabularnewline
evaporation mconstant $f$  & exponent $m$ in the concentration-dependent viscosity law\tabularnewline
evaporation tconstant $f$  & exponent in the relaxation-time law; use $f=1$ for Yarin et al. (2001)\tabularnewline
\uline{elastic modulus} $f$  & elastic modulus of the nanofiber\tabularnewline
external potential $f$  & electric potential between the nozzle and collector\tabularnewline
 & (the resulting electric field is oriented along \textit{x-axis})\tabularnewline
external potential freq $f$  & frequency of the waveform describing the electric potential\tabularnewline
external potential phase $f$  & initial phase of the waveform describing the electric potential\tabularnewline
external potential type $i$  & set the waweform describing the time evolution of the external\tabularnewline
 & electric potential. The \textit{integer} can be equal to 0 for a static
field,\tabularnewline
 & 1 for a rectangular waveform, \tabularnewline
 & 2 for a rectangular waveform of a RC circuit,\tabularnewline
 & 3 for orthogonal rotating electric field (default: 0)\tabularnewline
external potential time $f$  & time constant (relaxation time) of the RC circuit\tabularnewline
external potential vector $3f$  & vector of the electric field for \textit{x y z} components\tabularnewline
\uline{final time} $f$  & set the end time of the simulation \tabularnewline
\uline{finish}  & close the input file (last data record)\tabularnewline
gravity yes  & active the inclusion of the gravity force in the model (default: no)\tabularnewline
hbfluid yes  & allow the set of parameters needed to define\tabularnewline
 & the Herschel\textendash Bulkley stress. Note that if it is set no,
JETSPIN \tabularnewline
 & will use the simple Maxwellian fluid model (default: no)\tabularnewline
hbfluid consistency $f$  & flow consistency index rescaled by $\mu$ (in $\mu$ unit) as defined
in \tabularnewline
 & Eq \ref{eq:rescaled-prefactor} (Note you should also set: hbfluid
yes) (default: 1)\tabularnewline
hbfluid index $f$  & flow behavior index (Note you should also set: hbfluid yes) (default:
1)\tabularnewline
\uline{initial length} $f$  & length of the nanofiber at the initial time\tabularnewline
\uline{integrator} $i$  & set the integration scheme. The \textit{integer} can be '1' for\tabularnewline
 & \quad{}Euler, '2' for Heun, '3' for Runge-Kutta, and '4'\tabularnewline
 & \quad{}for the Platen scheme (only if airdrag force is activated).\tabularnewline
inserting yes  & inject new beads as explained in Subsec \tabularnewline
job time $f$  & Set job time to $f$ seconds\tabularnewline
kvfluid yes  & activate the Kelvin\textendash Voigt model which is used instead of
Maxwell\tabularnewline
 & model (default: no) (Note it is actually implemented only for working\tabularnewline
 & with the 3d-model without airdrag ativated: system 3 is mandatory)\tabularnewline
maximum displ $f$ & set the maximum displacement of the neighbor list (see Sec \ref{sec:Multiple-Timestep-Algorithm})\tabularnewline
multiple step yes & activate the multiple timestep algorithm for Coulomb forces (default:
no)\tabularnewline
multiple step every $i$ & set the interval of multiple-steps (see Sec \ref{sec:Multiple-Timestep-Algorithm})\tabularnewline
noise yes & active the inclusion of a random noise in the model (default: no)\tabularnewline
 & (note that if yes the directive ``integrator 4'' is mandatory)\tabularnewline
noise variance $f$ & variance of random noise in velocity space (default: 1 $\text{cm}^{2}/\text{s}^{2}$)\tabularnewline
noise diffusivity $f$ & diffusivity of random noise in velocity space (default: 0 $\text{cm}^{2}/\text{s}^{3}$)\tabularnewline
\uline{nozzle cross} $f$  & cross section radius of the nanofiber at the nozzle\tabularnewline
nozzle stress $f$  & viscoelastic stress of the nanofiber at the nozzle\tabularnewline
nozzle velocity $f$  & bulk fluid velocity of the nanofiber at the nozzle\tabularnewline
perturb yes  & active the periodic perturbation at the nozzle\tabularnewline
perturb freq $f$  & frequency of the perturbation at the nozzle\tabularnewline
perturb ampl $f$  & amplitude of the perturbation at the nozzle\tabularnewline
points $i$  & number of segment in which the jet is discretised\tabularnewline
primary cutoff $f$ & set the primary cutoff of the neighbor list (see Sec \ref{sec:Multiple-Timestep-Algorithm})\tabularnewline
print list $s_{1}\,\ldots$  & print on terminal statistical data as indicated by symbolic\tabularnewline
 & \quad{}strings (see Tab \ref{tab:outputlist} for detailed informations)\tabularnewline
print time $f$  & print data on terminal and output file every $f$ seconds\tabularnewline
print pdb frame $f$  & print the jet geometry in a single PDB file format in centimeters\tabularnewline
print pdb rescalepdb $f$  & print the PDB file format data rescaled by $f$ (default: 1)\tabularnewline
print xyz $f$  & print the trajectory in XYZ file format in centimeters\tabularnewline
 & \quad{}every $f$ seconds \tabularnewline
print xyz frame $f$  & print the jet geometry in a single XYZ file format in centimeters\tabularnewline
 & \quad{}every $f$ seconds (default name of files: frame\%06d.xyz)\tabularnewline
print xyz maxnum $i$  & print the XYZ file format with $i$ beads starting from the\tabularnewline
 & \quad{}nearest bead to the collector (default: 100)\tabularnewline
print xyz rescale $f$  & print the XYZ file format data rescaled by $f$ (default: 1)\tabularnewline
printstat list $s_{1}\,\ldots$  & print on output file statistical data as indicated by symbolic\tabularnewline
 & \quad{}strings (see Tab \ref{tab:outputlist} for detailed informations)\tabularnewline
removing yes  & remove beads at collector\tabularnewline
restart yes  & restart JETSPIN from a previous job (see Section \ref{sec:Restarting-JETSPIN})\tabularnewline
restart dump $i$  & print the \texttt{save.dat} file every $i$ steps\tabularnewline
restart reset  & reset all the accumulated statistical data collected in the previous\tabularnewline
 & simulation when a job is restarted\tabularnewline
resolution $f$  & discretization step length of nanofiber\tabularnewline
seed $i$  & seed control to the internal random number generator\tabularnewline
surface tension $f$  & surface tension coefficient of the nanofiber\tabularnewline
\uline{system} $i$  & set the nanofiber model. The \textit{integer} can be equal to \tabularnewline
 & \quad{}'1' for select the 1d-model, '3' for the 3d-model, and\tabularnewline
 &  '4' for the 3d-model with airdrag ativated\tabularnewline
\uline{timestep} $f$  & set the time step for the integration scheme\tabularnewline
\uline{viscosity} $f$  & viscosity of the nanofiber\tabularnewline
yield stress $f$  & yield stress (default: 0)\tabularnewline
\hline
\label{tab:inputlist}  & \tabularnewline
\end{longtable}

% Evaporation directives, included in the JETSPIN Data Files chapter.

\subsection{Evaporation directives\label{subsec:Evaporation-directives}}

The evaporation directives listed in Table~\ref{tab:inputlist} activate
and parameterize the model described in Section~\ref{sec:Evaporation}:
\begin{description}
 \item[\texttt{evaporation yes}] activate solvent evaporation and the
 concentration-dependent rheology.
 \item[\texttt{evaporation polymer frac} $f$] initial polymer mass
 fraction $c_{p,0}$.
 \item[\texttt{evaporation airviscosity} $f$] kinematic viscosity of
 air $\nu_a$ in $\mathrm{cm^2/s}$.
 \item[\texttt{evaporation temperature} $f$] temperature in kelvin.
 \item[\texttt{evaporation umidity} $f$] relative humidity $RH$ in the
 interval $[0,1]$. The historical spelling \texttt{umidity} is retained
 for input-file compatibility.
 \item[\texttt{evaporation diffusivity} $f$] vapour diffusion coefficient
 $D_a$ in $\mathrm{cm^2/s}$. If omitted,
 Eq.~\ref{eq:evap-default-diffusivity} is used at one atmosphere.
 \item[\texttt{evaporation relumidity} $f$] optional value of the
 saturation-vapour ratio $c_{s,\mathrm{eq}}(T)$. If omitted, it is
 evaluated from \texttt{evaporation temperature}. The historical keyword
 is retained for input-file compatibility.
 \item[\texttt{evaporation bconstant} $f$] parameter $B$ in
 Eq.~\ref{eq:evap-viscosity}.
 \item[\texttt{evaporation mconstant} $f$] exponent $m$ in
 Eq.~\ref{eq:evap-viscosity}.
 \item[\texttt{evaporation tconstant} $f$] exponent
 $t_{\mathrm{ev}}$ in the generalized relaxation-time law; use $f=1$
 to reproduce Yarin et al. (2001).
\end{description}

